\section{Cronograma}
\begin{table}[htbp]
\centering
\caption{Cronograma de atividades do semestre.}
\label{tab:cronograma_atividades}
\begin{tabular}{|p{3.5cm}|p{3.5cm}|p{7cm}|}
\hline
\textbf{Ponto de Controle} & \textbf{Data} & \textbf{Produto a ser entregue} \\ \hline
1º Ponto de Controle & Primeiro Mês & Primeira versão do artigo (Introdução, Referencial Teórico, Metodologia, Descrição das Implementações). Alinhamento com o orientador, possíveis correções antes da primeira entrega oficial. \\ \hline
2º Ponto de Controle & Segundo Mês & Segunda versão do artigo (Primeira versão revisada e acrescida de implementações, testes e resultados preliminares da coleta via GTFS-RT). \\ \hline
3º Ponto de Controle & Terceiro Mês & Terceira versão do artigo (Segunda versão revisada e acrescida dos resultados e análises via Modelos Mistos e Interpolação Linear). Essa margem é essencial para corrigir rotas caso algum teste ou implementação tenha falhado. \\ \hline
4º Ponto de Controle & Quarto Mês & Quarta Versão (Artigo completo e revisado pelo orientador). Deve conter: Resumo, Introdução, Revisão de literatura, Metodologia, Resultados (incluindo análise comparativa semanal), Conclusões e Trabalhos Futuros. Será a entrega oficial para revisão e liberação do orientador que será entregue a banca. \\ \hline
Bancas de Defesa & Quarto Mês (Data Oficial a ser marcada) & Apresentação oral para a banca. A avaliação será feita conjuntamente com o texto do artigo entregue. \\ \hline
5º Ponto de Controle & Até 48h a 5 dias após a banca & Quinta versão (Artigo final aprovado pelo orientador, contendo todas as sugestões da banca). Deve ser postado no Canvas e no site da Biblioteca. A nota da banca só é liberada após este envio. Se não postado, a nota é zerada. \\ \hline
\end{tabular}
\end{table}